% Resume: Arumalla Mohit Krishna
% Compile with pdfLaTeX.

\documentclass[letterpaper,10pt]{article}

\usepackage[margin=0.5in]{geometry}
\usepackage[T1]{fontenc}
\usepackage{lmodern}
\usepackage{microtype}
\usepackage{titlesec}
\usepackage{enumitem}
\usepackage[hidelinks]{hyperref}
\usepackage{tabularx}
\input{glyphtounicode}

\pdfgentounicode=1
\pagestyle{empty}
\urlstyle{same}
\raggedright
\setlength{\parindent}{0pt}
\setlength{\tabcolsep}{0pt}

\titleformat{\section}
  {\large\scshape}{}{0em}{}[\titlerule]

\titlespacing*{\section}{0pt}{6pt}{3pt}

\newcommand{\entry}[2]{%
  \noindent
  \begin{tabularx}{\linewidth}
    {@{}>{\raggedright\arraybackslash}X@{\hspace{6pt}}r@{}}
    \textbf{#1} & #2
  \end{tabularx}\par
}

\newcommand{\bullets}{
  \begin{itemize}[
    leftmargin=1.2em,
    topsep=2pt,
    itemsep=1pt,
    parsep=0pt,
    partopsep=0pt
  ]
}

\newcommand{\finishbullets}{
  \end{itemize}
}

\begin{document}

%----------HEADING----------
\begin{center}
    {\LARGE\bfseries Arumalla Mohit Krishna}\\[3pt]
    +1 (631) 352-8643 $|$
    \href{mailto:mohitkrishna.arumalla@stonybrook.edu}
    {mohitkrishna.arumalla@stonybrook.edu} $|$
    \href{https://linkedin.com/in/mohit-krishna-arumalla}
    {LinkedIn}
\end{center}
\vspace{-8pt}

%-----------EDUCATION-----------
\section{Education}

\entry
  {\textbf{Stony Brook University}}
  {Aug 2025 -- Dec 2026}

M.S. in Computer Science \hfill New York, USA\\
\textit{
Coursework: \textbf{Natural Language Processing}, \textbf{Introduction to Computer Vision}, Vision Language Modeling
}
\vspace{3pt}

\entry
  {\textbf{Birla Institute of Technology and Science, Pilani}}
  {Aug 2019 -- 2023}

B.E. in Electronics \& Communication; GPA: 8.02/10 \hfill India

%-----------PROFESSIONAL SUMMARY-----------
\section{Professional Summary}

Software engineer with 2+ years of experience in \textbf{developing, training, retraining, and deploying AI/ML models} addressing real-time business challenges, currently pursuing M.S. in Computer Science. Skilled in \textbf{selecting data sets, performing statistical analysis, running machine learning algorithms, and improving model performance}. Expertise in Python, PyTorch, cloud infrastructure, and scalable data pipelines supporting \textbf{efficient AI/ML solutions}. Strong understanding of model architecture, data integration, and performance metrics interpretation.

%-----------TECHNICAL SKILLS-----------
\section{Technical Skills}

\textbf{Languages:}
Python, Go, C++, Java, SQL, R\\

\textbf{ML \& AI Frameworks:}
PyTorch, TensorFlow/Keras, scikit-learn, NumPy, pandas,
Transformers, Vision Transformers, Self-Supervised Learning, LoRA, Visual Prompt Tuning,
Text Embeddings, Vector Search\\

\textbf{Data Engineering \& Systems:}
gRPC, Kafka, Apache Flink, Flink SQL, Debezium/CDC,
PostgreSQL, Redis, Elasticsearch, Airflow, BigQuery\\

\textbf{Infrastructure \& DevOps:}
NVIDIA Triton, AWS, Docker, Kubernetes, Linux, Git,
CI/CD, Istio/Envoy, New Relic

%-----------EXPERIENCE-----------
\section{Experience}

\entry
  {\textbf{Software Engineer II} $|$ \textbf{Infra.Market}, Bangalore}
  {Apr 2024 -- Aug 2025}

\bullets

  \item Designed and deployed \textbf{machine learning-powered semantic product search} with Amazon Titan text embeddings and Elasticsearch vector search, enhancing natural-language interfaces for product discovery.

  \item Engineered a CDC pipeline integrating Flink, Kafka, Debezium, PostgreSQL, RocksDB, and Elasticsearch; reduced synchronization latency from 10--15 minutes to 500 ms with exactly-once processing semantics.

  \item Deployed Istio/Envoy across 15+ microservices and integrated New Relic distributed tracing for production monitoring, supporting scalable and robust AI/ML infrastructure.

\finishbullets

\entry
  {\textbf{Software Engineer I} $|$ \textbf{Infra.Market}, Chennai}
  {Jul 2023 -- Mar 2024}

\bullets

  \item Owned Search and Location microservices; built gRPC services handling 500K+ daily requests to support real-time geolocation and product discovery with high availability.

  \item Reduced PostgreSQL connection leaks by 85\% through ORM tuning and fixed connection-pool exhaustion; validated improvements with load tests simulating 10K requests/second.

  \item Developed Airflow pipelines for automated data migration from BigQuery to Elasticsearch with query optimization and monitoring, ensuring data freshness and reliability for AI model training.

\finishbullets

%-----------PROJECTS-----------
\section{Projects}

\entry
  {GPU-Optimized Personalized Product Search}
  {Aug 2026 -- Oct 2026}

\textit{Python, PyTorch, BGE, NVIDIA Triton, Go, gRPC}

\bullets

  \item Developed \textbf{two-tower retrieval architecture} with \texttt{BAAI/bge-small-en-v1.5} embeddings and trained user tower combining query embeddings with interaction history to generate personalized recommendations.

  \item Built inference pipeline with precomputed product embeddings and approximate nearest-neighbor retrieval via a Go/gRPC gateway; integrated NVIDIA Triton for efficient GPU query-encoder serving.

  \item Designed performance benchmarks to measure latency, throughput, GPU utilization, and cost; evaluated ranking quality using Recall@K and NDCG@K, implementing recovery strategies for overload scenarios.

\finishbullets

\entry
  {Attention-Based Hallucination Detection in NMT}
  {Apr 2026 -- May 2026}

\textit{Python, PyTorch, mBART, Transformers, scikit-learn}

\bullets

  \item Built an English-German hallucination detection pipeline leveraging teacher-forced mBART attention features from WMT22/WMT23 quality estimation datasets.

  \item Extracted and analyzed cross- and self-attention statistics across multiple decoder layers and heads, including entropy and attention mass ratios.

  \item Trained token-, sentence-, and head-level classifiers achieving F1 score of 0.712 and ROC-AUC of 0.695; proposed hallucination-aware reranking leveraging attention-based risk scores.

\finishbullets

\entry
  {Prompt/LoRA for Self-Supervised Learning}
  {Ongoing}

\textit{Python, PyTorch, Vision Transformers}

\bullets

  \item Investigating domain-specific visual prompting and \textbf{parameter-efficient LoRA adaptation} in a self-supervised ViT framework to learn transferable image representations.

  \item Comparing backbone performance with and without prompts/LoRA at inference; analyzing the impact on classification and robustness under domain shift.

\finishbullets

\end{document}